\documentclass[11pt,a4paper]{article}
\usepackage[margin=25mm]{geometry}
\usepackage{amsmath,amssymb,amsthm}
\newtheorem{proposition}{Proposition}
\newtheorem{corollary}[proposition]{Corollary}
\newcommand{\Aut}{\operatorname{Aut}}
\newcommand{\Inn}{\operatorname{Inn}}
\newcommand{\Out}{\operatorname{Out}}
\title{Counting chiral markings by stabilizers}
\author{Human Compression project}\date{}
\begin{document}\maketitle
The second atlas in Hartley, Hubard and Leemans,
\emph{Two atlases of abstract chiral polytopes for small groups}
(2012), specifies many finite automorphism groups. We apply the already
reviewed prefix characterization and enumeration theorem to each group.
This note supplies an orbit count and its concrete implementation.
No numerical entry of the second atlas is assumed as an input.

For a finite group $G$, an admissible prefix is an ordered tuple
$(1,a_1,\ldots,a_{n-1})$ satisfying the acquired adjacent-order,
nonadjacent-involution and difference-subgroup intersection conditions.
For $J\subseteq\{0,\ldots,n-1\}$ write
\[
D_J=\langle a_i^{-1}a_j:i,j\in J\rangle.
\]
The shared coordinate theorem identifies these with the corresponding
interval-rotation subgroups. The shared search theorem proves that
prefixes remain admissible, each extension strictly enlarges the
generated subgroup, and an admissible generating tuple cannot be
extended further. The Schulte--Weiss reconstruction theorem remains
the attributed geometric input.

\begin{proposition}
Choose one representative $a_1$ from each conjugacy class of elements
of order greater than two in $G$. At a proper prefix generating $H$,
choose one representative of each $C_G(H)$-orbit on the permissible
next involutions, and retain precisely the extensions passing the
intersection tests. Stop a branch when its entries generate $G$.
The generating tuples so obtained are representatives of exactly all
inner automorphism orbits of admissible generating tuples.
\end{proposition}
\begin{proof}
Conjugation preserves every admissibility condition. After the first
entry has been moved to its chosen representative, precisely its
centralizer remains available to move the next entry. More generally,
an element fixes every entry of a prefix if and only if it centralizes
the subgroup they generate. The permissible next entries form an
invariant set under this centralizer. Choosing one representative of
each such orbit therefore gives exactly one representative of each
orbit of the extended prefixes over the fixed prefix. Induction proves
both completeness and uniqueness of the orbit representatives at every
depth. The acquired hereditary and strict-growth statements justify
the retention tests and stopping rule in every possible rank.
\end{proof}

The implementation first filters the involutions by the order
conditions and by exclusion from $H$. These tests are invariant under
$C_G(H)$, so they can precede the orbit computation. The intersection
property is also invariant; it is tested on one representative per
orbit. Every difference subgroup is computed from a star at its least
index, by the shared star-generation identity. The checker tests all
pairs of subsets, including empty intersections. It does not substitute
a necessary condition for the complete intersection property.

\begin{proposition}
Let $\mathcal T$ be a finite, $\Aut(G)$-invariant set of ordered
generating tuples. If $b$ is the number of its inner automorphism orbits,
then the number of its full automorphism orbits is
\[
\frac{b}{|\Out(G)|},\qquad
|\Out(G)|=\frac{|\Aut(G)|\,|Z(G)|}{|G|}.
\]
\end{proposition}
\begin{proof}
An automorphism fixing all entries of a generating tuple fixes every
element of $G$. Thus $\Aut(G)$ acts freely on $\mathcal T$, and so does
its subgroup $\Inn(G)$. Each full orbit has $|\Aut(G)|$ elements and
each inner orbit has $|\Inn(G)|=|G|/|Z(G)|$ elements. A full orbit is
a union of precisely $|\Out(G)|$ inner orbits. Dividing the total
number of inner orbits by this common number proves the claim.
\end{proof}

Rank, the ordered Schlaefli type, and the existence of a reflection
automorphism are invariant under $\Aut(G)$. The proposition therefore
applies separately in each such class. The reflection test is the
acquired criterion that an automorphism send every $a_i$ to $a_i^{-1}$.
A homomorphism with these images, if it exists, is surjective because
the images generate $G$; finiteness makes it an automorphism. We use
the checked homomorphism constructor, not its unchecked variant.
An unsuccessful constructor excludes the reflection relative to the
stated exact group-arithmetic premise.

Consequently the number of chiral oriented classes of a given rank and
type is its number of nonreflectable inner representatives divided by
$|\Out(G)|$. The acquired mirror theorem pairs these classes freely
when orientation is forgotten. Duality is still kept distinct.
The regular tuples in this calculation concern rotation groups;
they are not identified with the source's separate column counting
regular polytopes by their full automorphism groups.

\begin{corollary}
A finite chiral polytope cannot have an adjacent Schlaefli entry two.
\end{corollary}
\begin{proof}
The corresponding flag adjacencies commute. Together with the distant
commuting relations, this partitions the colours into two nonempty
sets $I,J$ whose generated permutation groups $M_I,M_J$ commute.
Let $W=M_I\times M_J$ act on the actual flags through multiplication
of these commuting permutations. It acts transitively, so the flags
are $N\backslash W$ for the stabilizer $N$ of a chosen flag.
The previously proved flag-component argument gives
$N=(N\cap M_I)\times(N\cap M_J)$: if $ab\in N$, the flag reached
by $a$ is also reached by $b^{-1}$, and these two paths preserve
complementary sets of faces, hence preserve the whole base flag.
The normalizer splits as a product as well. Chirality would put every
even word $r_ir_j$ in this normalizer. Taking $i\in I,j\in J$ then
puts both individual generators in their respective normalizer
factors. All generators normalize $N$, making the polytope regular,
a contradiction. This uses the actual flag action; no unproved
intersection property is imposed on the monodromy group.
\end{proof}

\paragraph{Verification boundary and attribution.}
The coordinate and prefix theory, reflection criterion and geometric
reconstruction are acquired from the reviewed manuscript on the
algorithm of Buekenhout, Leemans and Tranchida, with the attributions retained
there. The centralizer recursion uses the same orbit argument
as other parts of the project. The count of free orbits is elementary;
no priority claim is made. The last corollary extends the previously
reviewed product proof by using the direct product of the two actual
monodromy subgroups.

The program retains GAP's exact group operations, including complete
conjugacy classes, centralizers, orbits, subgroup intersections,
automorphism groups and homomorphism construction, as software premises.
Literal small-group controls compare its complete output with the
independently acquired Cayley-table search. Large finite target calls
require their own complete receipts, input identities and distinct
review; this note alone confers no source census or whole-paper credit.
\end{document}
